\documentclass[11pt, a4paper]{article}
\usepackage[margin=1in]{geometry}
\usepackage{hyperref}

\begin{document}

\noindent
\begin{tabular}{@{}l}
Maciej Tadej \\
Mathematical Institute, University of Wrocław \\
pl. Grunwaldzki 2, 50-384 Wrocław, Poland \\
\href{mailto:maciej.tadej@math.uni.wroc.pl}{maciej.tadej@math.uni.wroc.pl}
\end{tabular}

\vspace{2em}

\noindent
To the Editorial Board \\
\textbf{Integral Equations and Operator Theory} \\
Springer

\vspace{2em}

\noindent
Dear Editors,

\vspace{1em}

\noindent
I am pleased to submit the manuscript entitled \textbf{"Eigenvalues and eigenfunctions of the non-local dispersal operator with symmetric kernel and Neumann-type boundary condition"} for consideration for publication in the \textit{Integral Equations and Operator Theory}.

This study solves an open problem in the spectral theory of non-local dispersal operators on bounded domains with continous, symmetric kernel. Unlike the classical Laplacian operator, these type of integral operators lack compactness, which makes the study of discrete eigenvalues very difficult. While the non-local Dirichlet problem is well understood already, finding non-trivial eigenfunctions for the non-local Neumann boundary condition on bounded domains remained unsolved.

In this paper, we provide a rigorous variational construction of a sequence of non-trivial eigenfunctions. The main novelties of our work are:
\begin{itemize}
    \item We define geometric conditions for existence of the spectral gap, linking the domain's size and shape with the kernel's variance.
    \item We construct a finite-dimensional Galerkin approximation and rigorously prove the strong $L^2$-convergence of the minimizing sequences to the exact eigenpairs.
\end{itemize}

These results are important for both theory and practical applications, particularly in ecology. Non-local operators are widely used to model real-world phenomena such as seed dispersal, vegetation pattern formation, and species survival in isolated habitats. By considering bounded domains with Neumann-type boundary conditions, our framework models habitats where individuals cannot cross the boundary, offering a much more realistic setting than artificial unbounded domains or periodic boundary conditions.

This manuscript has not been published previously and is not under consideration for publication elsewhere. All authors approve its submission.

I would like to suggest the following Handling Editors whose research focus closely aligns with the spectral theory and variational methods used in this manuscript:
\begin{itemize}
    \item \textbf{Prof. Jussi Behrndt} (Graz University of Technology).
    \item \textbf{Prof. Tom ter Elst} (University of Auckland).
\end{itemize}

Thank you for your time and consideration.

\vspace{1em}

\begin{flushright}
Sincerely, \\
Maciej Tadej
\end{flushright}

\end{document}